\documentclass[a4paper]{article}
\usepackage[UTF8,fontset=none]{ctex}
\usepackage{graphicx}
\usepackage{caption}
\usepackage{array}
\usepackage{fancyhdr}
\usepackage{lastpage}
\usepackage{geometry}
\usepackage{fontspec}
\usepackage{amsmath}
\usepackage{amsfonts}
\usepackage{ulem}
\usepackage{gensymb}
\usepackage{multirow}
\IfFileExists{siunitx.sty}{\usepackage{siunitx}}{}
\usepackage{etoolbox}

% 本文件由 acxy5hfo.tex 派生，仅用于模板兼容性和一般排版测试。
% 所有姓名、学号和实验数据均为醒目标注的虚构示例，不可作为实验记录提交。

\makeatletter
\setlength{\textfloatsep}{5pt plus 2pt minus 4pt}
\setlength{\floatsep}{5pt plus 2pt minus 4pt}
\setlength{\intextsep}{5pt plus 2pt minus 4pt}
\makeatother

\setlength{\parindent}{0pt}

% 原模板硬编码 Windows 字体；本机测试副本改用已安装的中西文字体。
\setmainfont{TeX Gyre Termes}
\setsansfont{TeX Gyre Heros}
\setCJKmainfont[AutoFakeBold=2]{AR PL SungtiL GB}
\setCJKsansfont[AutoFakeBold=2]{Droid Sans Fallback}

\geometry{left=2cm,right=2cm,top=2.5cm,bottom=2.5cm}

\pagestyle{fancy}
\fancyhf{}
\renewcommand{\headrulewidth}{0pt}
\fancyfoot[C]{浙江大学物理实验教学中心制 \quad 第 \thepage 页 / 共 \pageref{LastPage} 页}

\begin{document}

\linespread{1.5}
\selectfont

\begin{center}
\vspace*{1em}
\includegraphics[width=6.2cm]{assets/zju-name.jpg}\\
\vspace{0.6cm}
{\fontsize{50pt}{54pt}\selectfont \textbf{物理学实验I}}\\
\vspace{2em}
{\fontsize{50pt}{54pt}\selectfont \textbf{电子实验报告}}\\
\vspace{10em}

\Large \textbf{实验名称：\uline{\makebox[6.5cm][c]{单摆周期测量（排版示例）}}}\\
\Large \textbf{实验桌号：\uline{\makebox[6.5cm][c]{A-00（示例）}}}\\
\Large \textbf{指导教师：\uline{\makebox[6.5cm][c]{示例教师}}}\\
\vspace{3em}
\Large \textbf{班 \quad 级：\uline{\makebox[6cm][c]{示例班级}}}\\
\Large \textbf{姓 \quad 名：\uline{\makebox[6cm][c]{示例学生}}}\\
\Large \textbf{学 \quad 号：\uline{\makebox[6cm][c]{0000000000}}}\\
\vspace{3em}
\Large \textbf{实验日期：2026年9月21日~~星期一下午（示例）}\\
\vspace{1em}
\normalsize （本报告全部内容仅用于模板排版测试）
\end{center}

\newpage

\section*{一、预习报告}

\subsection*{1. 实验综述}
本示例以小角度单摆为对象，使用秒表记录多次摆动的总时间，并由
$T=t/N$ 得到单次周期。理想条件下，周期与摆长满足
\begin{equation}
  T=2\pi\sqrt{\frac{L}{g}}, \qquad
  g=\frac{4\pi^2L}{T^2}.
\end{equation}
通过重复测量、求平均值和估计不确定度，可检验理论关系以及人工计时对结果的影响。

\subsection*{2. 实验重点}
规范读取摆长和累计计时数据，使用合适的有效数字表达平均值与不确定度，并保证表格、公式和单位清楚一致。

\subsection*{3. 实验难点}
减小释放角度与人为启停秒表造成的系统偏差，并区分随机不确定度和仪器分辨率带来的影响。

\section*{二、原始数据}
下框模拟插入“教师签名的原始记录照片”后所占的版面；真实报告中应替换为清晰、完整且未经裁掉签字信息的扫描图。

\begin{figure}[htbp]
  \centering
  \fbox{\parbox[c][5.6cm][c]{0.88\textwidth}{\centering
  原始数据记录照片占位区域\\[0.8em]
  （排版示例：此处没有真实实验图片或教师签名）}}
  \caption*{图像占位框仅用于检查宽度、分页和图文间距}
\end{figure}

% 避免下一节标题与其首个结果表跨页分离。
\clearpage
\section*{三、结果与分析}
\subsection*{1. 数据处理与结果}
表中数值为专门构造的排版示例，不代表任何一次实际实验。

\begin{table}[htbp]
  \centering
  \caption{单摆周期测量的示例数据}
  \begin{tabular}{|c|c|c|c|}
    \hline
    次数 & 摆动次数 $N$ & 总时间 $t$/s & 周期 $T$/s \\
    \hline
    1 & 20 & 40.12 & 2.006 \\
    2 & 20 & 40.08 & 2.004 \\
    3 & 20 & 40.16 & 2.008 \\
    4 & 20 & 40.10 & 2.005 \\
    5 & 20 & 40.14 & 2.007 \\
    \hline
    \multicolumn{3}{|c|}{平均值} & 2.006 \\
    \hline
  \end{tabular}
\end{table}

取示例摆长 $L=(0.9980\pm0.0005)\,\mathrm{m}$，平均周期
$\overline{T}=2.006\,\mathrm{s}$，代入式（1）得到
\[
  g_{\mathrm{example}}=\frac{4\pi^2\times0.9980}{(2.006)^2}
  =9.79\,\mathrm{m\,s^{-2}}.
\]
若周期和摆长相互独立，相对合成不确定度可写为
\[
  \frac{u(g)}{g}=\sqrt{\left(\frac{u(L)}{L}\right)^2+
  \left(2\frac{u(T)}{T}\right)^2}.
\]

\subsection*{2. 误差分析}
示例数据的离散程度较小，但人工计时的反应延迟可能在每次启停时引入偏差。累计测量多个周期能够降低单次反应时间在总时间中的相对占比。若摆角不够小、摆长测量未取到摆球质心，或支点存在摩擦，结果还会产生系统偏差。

\subsection*{3. 实验探讨}
一般报告应同时给出原始量、处理方法、最终结果与不确定度，避免只列计算结果。表格标题、变量符号和单位应能让读者脱离正文仍理解各列含义。

\section*{四、思考题}
\textbf{题目：为什么测量多个完整周期，而不直接测量一个周期？}

\textbf{答：}在反应时间近似不变的情况下，累计时间随周期数增加，而启停误差不会同比例增加，因此换算到单个周期后的相对误差通常更小。周期数也不宜无限增加，否则阻尼、振幅变化和计数失误可能变得明显。

\vfill
\begin{center}
  \small —— 一般用例结束；所有字段与数据均为虚构示例 ——
\end{center}

\end{document}
